\documentclass{article}
\usepackage{enumitem}
\usepackage{amsmath}
\begin{document}

\title{Chapter 1: The Living World}
\date{}
\maketitle

\begin{enumerate}[label=Q\arabic*.]

\item Which of the following best defines metabolism in the context of living organisms?
\\(A) Breakdown of food molecules only
\\(B) Sum total of all chemical reactions occurring in a living body
\\(C) Synthesis of proteins from amino acids
\\(D) Exchange of gases with the environment

\item A scientist discovers an object that grows, reproduces, but shows no cellular organization and no metabolism. This object is most likely:
\\(A) A virus
\\(B) A prion
\\(C) A crystal
\\(D) A dormant seed

\item Which of the following is NOT a defining property of living organisms according to NCERT?
\\(A) Growth
\\(B) Reproduction
\\(C) Locomotion
\\(D) Cellular organization

\item Taxonomic hierarchy in correct descending order is:
\\(A) Kingdom $\rightarrow$ Phylum $\rightarrow$ Class $\rightarrow$ Order $\rightarrow$ Family $\rightarrow$ Genus $\rightarrow$ Species
\\(B) Kingdom $\rightarrow$ Class $\rightarrow$ Phylum $\rightarrow$ Order $\rightarrow$ Family $\rightarrow$ Genus $\rightarrow$ Species
\\(C) Phylum $\rightarrow$ Kingdom $\rightarrow$ Class $\rightarrow$ Order $\rightarrow$ Family $\rightarrow$ Species $\rightarrow$ Genus
\\(D) Kingdom $\rightarrow$ Phylum $\rightarrow$ Order $\rightarrow$ Class $\rightarrow$ Family $\rightarrow$ Genus $\rightarrow$ Species

\item The binomial nomenclature system was introduced by:
\\(A) Charles Darwin
\\(B) Carolus Linnaeus
\\(C) Ernst Mayr
\\(D) Gregor Mendel

\item In binomial nomenclature, which of the following is correctly written?
\\(A) \textit{homo sapiens}
\\(B) \textit{Homo Sapiens}
\\(C) \textit{Homo sapiens}
\\(D) \textit{HOMO SAPIENS}

\item Which taxonomic category contains the most closely related organisms?
\\(A) Order
\\(B) Family
\\(C) Genus
\\(D) Species

\item A taxon is best defined as:
\\(A) A rank in the taxonomic hierarchy
\\(B) A unit of classification representing any level of hierarchy
\\(C) The species name in binomial nomenclature
\\(D) The study of identifying organisms

\item The correct scientific name of the common housefly is:
\\(A) \textit{Musca domestica}
\\(B) \textit{Apis mellifera}
\\(C) \textit{Felis catus}
\\(D) \textit{Mangifera indica}

\item Which of the following properties distinguishes living from non-living most precisely?
\\(A) Growth
\\(B) Response to stimuli
\\(C) Consciousness and self-awareness
\\(D) Ability to reproduce using genetic material

\item ICZN stands for which regulatory body in taxonomy?
\\(A) International Code of Zoological Nomenclature
\\(B) International Commission of Zoological Names
\\(C) International Classification of Zoological Nomenclature
\\(D) Indian Code of Zoological Nomenclature

\item Which of the following represents a correctly identified specimen in a herbarium?
\\(A) Dried plant pasted on sheet with locality, date, and collector's name
\\(B) A living plant in a greenhouse with a name tag
\\(C) A photograph of a plant with common name
\\(D) A plant specimen preserved in formalin

\item Assertion: Viruses are considered at the border of living and non-living.\\
Reason: Viruses can replicate only inside a host cell using host machinery.
\\(A) Both true, Reason is correct explanation
\\(B) Both true, Reason is not correct explanation
\\(C) Assertion true, Reason false
\\(D) Assertion false, Reason true

\item Which of the following best represents the concept of ``species'' according to the biological species concept?
\\(A) Organisms with identical morphology
\\(B) Organisms that share a common ancestor
\\(C) Organisms that can interbreed and produce fertile offspring
\\(D) Organisms occupying the same ecological niche

\item In taxonomy, the term ``type specimen'' refers to:
\\(A) The most common organism of a species
\\(B) The single specimen on which the original species description is based
\\(C) The largest specimen in a museum collection
\\(D) A specimen used for genetic analysis

\end{enumerate}
\end{document}
